\documentclass[10pt]{beamer}

\usetheme{Madrid}
\usecolortheme{default}
\setbeamertemplate{navigation symbols}{}
\setbeamertemplate{footline}{\hfill\insertframenumber/\inserttotalframenumber\hspace*{5pt}\vspace{2pt}}

\title{Project Summary}
\subtitle{Adaptive EEG-to-Spike Framework}
\author{
  Bala Mosay J \quad (2023BTIA029) \\[2pt]
  Allwyn Jeffo Raj A \quad (2023BTIA002)
}
\institute{Under the guidance of Dr.\ Vaibhav Meshram}
\date{September 2, 2026}

\begin{document}

\begin{frame}
\titlepage
\end{frame}

\begin{frame}{What We Coded: Data Pipeline}

\small
\textbf{5 files handle all data loading and cleaning:}

\vspace{4pt}
\begin{itemize}
  \item \texttt{loader.py} --- Downloads brain data from PhysioNet for 103 subjects
  \item \texttt{preprocessing.py} --- Cleans signals: bandpass filter, ICA, epoching, AutoReject
  \item \texttt{events.py} --- Labels each trial (left hand, right hand, feet, rest)
  \item \texttt{balancer.py} --- Makes sure all classes have equal trial counts
  \item \texttt{normalization.py} --- Scales data per subject (z-score, feature scaling)
\end{itemize}

\vspace{4pt}
All 103 subjects' cleaned data is cached so preprocessing only runs once.

\end{frame}

\begin{frame}{What We Coded: Features + SNN}

\small
\textbf{7 files build the feature extractor and the brain-like AI:}

\vspace{4pt}
\begin{itemize}
  \item \texttt{rcsp.py} --- Regularized CSP spatial filter (finds best brain patterns)
  \item \texttt{bandpass.py} --- Applies CSP per frequency band, outputs 32 features per trial
  \item \texttt{encoder.py} --- Converts numbers to spikes (binary 0/1 using threshold)
  \item \texttt{lif\_net.py} --- The SNN: Linear, LIF, Dropout, Linear, LIF, Linear (64 hidden)
  \item \texttt{trainer.py} --- Trains the network: Adam optimizer, cosine LR, gradient clipping
  \item \texttt{evaluator.py} --- Measures accuracy, kappa, F1 score
  \item \texttt{synops.py} --- Counts SynOps (how many neuron connections fire = energy usage)
\end{itemize}

\end{frame}

\begin{frame}{What We Coded: Experiments + Tests}

\small
\textbf{3 experiment runners:}

\vspace{4pt}
\begin{itemize}
  \item \texttt{run\_baseline.py} --- Runs full pipeline on all 103 subjects (5-fold CV)
  \item \texttt{run\_logreg\_diagnostic.py} --- Tests simple classifier to find bottleneck
  \item \texttt{run\_per\_subject\_diagnostic.py} --- Tests each subject individually
\end{itemize}

\vspace{6pt}
\textbf{39 tests, all passing:}

\vspace{4pt}
\begin{itemize}
  \item Tests for every module: no data leakage, shapes correct, encoding works
  \item Leakage guard: label-dependent steps are fit on training folds
  \item ICA/AutoReject run per subject before CV (transductive, label-free)
\end{itemize}

\end{frame}

\begin{frame}{Results So Far}

\small
\textbf{Subject-mixed trial-level evaluation (103 subjects, 5-fold CV):}

\vspace{6pt}
\begin{itemize}
  \item Chance (random guess): 25.0\%
  \item Our SNN (baseline): \textbf{31.8\%} +/- 1.1
  \item LogReg (same features): \textbf{34.6\%} +/- 0.8
\end{itemize}

\vspace{6pt}
\textbf{Key finding:}

\vspace{4pt}
\begin{itemize}
  \item The SNN (31.8\%) is only 2.8\% behind LogReg (34.6\%) on the same features
  \item The model is not the bottleneck --- the features are
  \item Within a single subject, LogReg achieves 42.0\% --- features do work
  \item Unseen-subject generalization is the challenge (LOSO pending)
\end{itemize}

\end{frame}

\begin{frame}{Phase 3: What's Next}

\small
\textbf{3 novel contributions to improve accuracy:}

\vspace{4pt}
\begin{enumerate}
  \item \textbf{Adaptive Band Selection (C1)} --- Pick best frequencies per person using KNN Mutual Information instead of fixed mu/beta bands
  \item \textbf{Learnable Spike Threshold (C2)} --- Make threshold a learnable parameter + add sparsity penalty to save energy
  \item \textbf{Dual-Branch PLIF Network (C3)} --- Two parallel branches (fast tau=5ms, slow tau=30ms) combined with attention mechanism
\end{enumerate}

\vspace{6pt}
\textbf{Target:}

\vspace{4pt}
\begin{itemize}
  \item Accuracy: 55--65\% (up from 31.8\%)
  \item Lower SynOps than baseline (more energy efficient)
\end{itemize}

\end{frame}

\end{document}
